% =============================================================================
% PARTE II: La Mente Fuori Posto
% =============================================================================

\addcontentsline{toc}{part}{PARTE II: La Mente Fuori Posto} % Questo la aggiunge all'indice

\cleardoublepage
\thispagestyle{empty}

\vspace*{\fill}

\begin{center}
    % Numero della parte
    {\sffamily\huge\textls[100]{\MakeUppercase{Parte II}}}
    
    \vspace{1.5em}
    
    % Linea decorativa
    \rule{0.5\textwidth}{0.4pt}
    
    \vspace{1.5em}
    
    % Titolo della parte
    {\sffamily\Huge\bfseries La Mente Fuori Posto}
    
    \vspace{2em}
    
    % Sottotitolo/descrizione
    \begin{minipage}{0.7\textwidth}
        \centering
        {\large\itshape\color{moderngray}
        I nostri bias in azione nel mondo moderno.\\
        Casi di studio di una struttura cerebrale antica\\
        in un mondo nuovo.}
    \end{minipage}
    
    \vspace{1.5em}
    
    % Linea decorativa
    \rule{0.5\textwidth}{0.4pt}
\end{center}

\vspace*{\fill}

\clearpage
\thispagestyle{empty}
\vspace*{2em}
\begin{center}
	{\sffamily\large\bfseries\color{darkgray} Dalla mente al mondo}
\end{center}
\vspace{1em}

Abbiamo visto che la mente non registra il mondo, lo ricostruisce. La memoria riscrive il passato per proteggere chi siamo, il pilota rapido decide prima che quello lento se ne accorga, la dissonanza trasforma le prove scomode in giustificazioni, e il gruppo può correggere questi errori oppure moltiplicarli. Finora, però, li abbiamo osservati quasi in laboratorio.

La domanda di questa parte è che cosa succede quando quella costruzione esce di casa. Che cosa accade quando decide di risparmi, di diagnosi, di sentenze, del clima in cui vivranno i nostri figli? E da dove vengono le cornici dentro cui decidiamo, le figure di cui ci fidiamo, le storie con cui siamo cresciuti, le idee su uomini e donne, su noi e gli altri?

La parte si muove in tre tempi. Il primo riguarda le decisioni e le loro conseguenze, in quattro campi dove un errore lascia tracce misurabili. Il secondo riguarda la costruzione culturale della realtà, cioè ciò che pensiamo prima ancora di pensare. Il terzo riguarda la trasmissione, il modo in cui le convinzioni passano da una generazione all'altra e si rinforzano guardando solo chi ce l'ha fatta.

Questo passaggio è necessario per una ragione precisa. Ogni decisione di cui parleremo lascia un documento, una cartella clinica, una sentenza, un prezzo, una storia raccontata. Quei documenti sono il passato da cui, nella Parte III, impareranno le macchine.

\vspace{1em}
\begin{storia}
Anna ha cinquantadue anni, un senso di peso al petto, nausea, il fiato corto. Nella sua cartella, da anni, c'è scritto «disturbo d'ansia». Giulia la conosce bene, la vede spesso, e quella diagnosi le viene incontro prima ancora che Anna finisca di parlare. Il ricordo che di solito la mette in allarme ha il volto di un ragazzo, non di una donna di mezza età già nota per l'ansia. Giulia scrive «sintomi ansiosi, rivalutare» e le raccomanda di riposare. La sera la figlia di Anna insiste per andare in pronto soccorso. È un infarto. Anna se la cava. La nota resta dov'è, una riga tra milioni di righe.
\end{storia}

\cleardoublepage
